\documentclass[11pt,letterpaper]{article}
\usepackage[margin=0.85in]{geometry}
\usepackage[T1]{fontenc}
\usepackage{lmodern,microtype,amsmath,amssymb,booktabs,longtable,array}
\usepackage{enumitem,fancyvrb,fancyhdr,xcolor}
\usepackage[colorlinks=true,linkcolor=black,urlcolor=blue!45!black]{hyperref}
\setlength{\parindent}{0pt}
\setlength{\parskip}{6pt}
\setlength{\emergencystretch}{3em}
\setlist{nosep,leftmargin=*}
\renewcommand{\arraystretch}{1.18}
\pagestyle{fancy}
\fancyhf{}
\fancyhead[L]{\small CLAS12 track finding}
\fancyhead[R]{\small Evaluation specification v1}
\fancyfoot[C]{\thepage}
\setlength{\headheight}{14pt}
\newcommand{\metric}[1]{\texttt{\small\detokenize{#1}}}
\newcommand{\code}[1]{\texttt{\detokenize{#1}}}
\newcommand{\NA}{\text{undefined}}
\newcommand{\ratio}[2]{\frac{#1}{#2}}
\newcommand{\file}[1]{\texttt{\small\path{#1}}}
\DefineVerbatimEnvironment{Code}{Verbatim}{fontsize=\small}

\begin{document}
\begin{center}
{\LARGE\bfseries Evaluating eligible MC track reconstruction}\par
\vspace{6pt}
{\large Metric definitions, provenance, aggregation and validation}\par
\vspace{5pt}
\code{mc_reconstruction_v1}\quad |\quad 10 November 2026\par
Phase 1: shared offline evaluator
\end{center}

\textbf{Benchmark objective.} Reconstruct the eligible simulated particle tracks
as completely and cleanly as possible. Clusters from other eligible tracks,
excluded MC particles and measured overlay reduce the score when they enter an
eligible track's associated prediction. Independently reconstructed predictions
containing only non-target clusters do not reduce the primary score. Their
physical validity is not established by the available target labels.

\textbf{Primary metric.} \metric{mc_track_dice_mean} is the sum of event-local,
one-to-one matched Dice scores divided by the number of eligible truth tracks
across the dataset. Unmatched truth tracks contribute zero. It is not a
track-count F1 score, an event average, or a training loss.

\textbf{Scope.} This specification defines an opt-in evaluation API and a saved-label
CPU scorer. The independent accepted matcher, loss formulas, model architecture,
checkpoint-selection rules and prediction thresholds are unchanged. A subsequent,
explicitly authorized downstream eligibility patch is described in Section~\ref{sec:cut};
it changes supervision without rewriting the skimmer product.

\textbf{Provenance boundary.} A non-target label does not establish measured
overlay. The overlay metrics in this version count \emph{verified pure-overlay
clusters}. Mixed and unresolved clusters are reported separately; the stored
scalar provenance does not support a complete member-level overlay fraction.

\tableofcontents
\newpage

\section{Eligibility is a fixed scientific policy}

The benchmark targets are primary MC tracks satisfying the established ownership
requirements and having at least four CVT clusters. Secondary particles, decay
products, and tracks with fewer than four qualifying clusters are not additional
targets. This policy is not reconsidered by the matching algorithm.

The evaluator consumes the existing segmentation labels. A valid cluster with a
non-background label belongs to the corresponding supplied target track. The
default background label is $-1$. It means \emph{not assigned to an eligible target
in this label contract}; it is not a statement that the cluster has no MC origin.
The evaluator does not convert an excluded MC particle into a target, revise
ownership, impose a new size cut, or alter any training labels.

\subsection{Ownership and eligibility are distinct from provenance}

In the inspected HDF5 producer, direct primary membership uses the simulation
condition $\mathrm{tid}=\mathrm{otid}$. The compact ownership rule requires an
unambiguous direct-primary contributor with strictly more than half of the
observed cluster members. Contested ownership is unassigned. The persisted
\code{primary_segment} is a compact label derived from the resulting positive
primary owner IDs. The raw source view preserves it; the two adapters now apply
the downstream eligibility projection in Section~\ref{sec:cut}.

Cluster provenance answers a different question: which observed members came
from MC, measured overlay, or unresolved sources? A primary-owned cluster can
contain other members. A non-target cluster can contain excluded MC. Neither
observation changes eligibility.

\subsection{Policy compliance must be audited without relabeling}

The HDF5 projection at the time of the initial Phase-1 audit did not apply a four-cluster track cut.
The locally available qualification HDF5 sample contains positive
\code{primary_segment} labels with fewer than four clusters: the first 1,000
events inspected contained 125 such labels among 1,900 labelled tracks. This is
evidence about that qualification sample, not a claim about every campaign input.

The shared scoring API preserves supplied metric inputs exactly. It records the
short-track count in \metric{n_truth_tracks_below_policy_min_clusters}.
It does not drop those tracks.
A nonzero count flags an input that cannot yet be described as satisfying the
four-cluster policy on the evaluated support. Such an input can be rescored for
historical parity, but its result must be described as conditional on its supplied
labels. The authorized downstream fix is separate from the metric computation.
Toy tests can deliberately supply short tracks to
verify this non-relabeling guarantee.

The scorer must see complete events. Cropping an eligible track can both shrink
its denominator and produce a short-track flag; it does not justify redefining
the eligibility rule. Filtering difficult events or excluding missed targets
after inference would change the benchmark population.

\subsection{Authorized downstream eligibility enforcement}\label{sec:cut}

The downstream label policy \code{primary_mc_min4_v1} now enforces the existing
four-cluster requirement for both CLAS12 track-finding and noise-tagging adapters.
For each complete event, the shared \code{eligible_primary_labels} projection
counts the existing nonnegative primary-segment labels and maps every set of
size one, two or three to $-1$. Sets of size four or more retain their original
IDs. No clusters, physical owner IDs or provenance fields are removed or rewritten.

This happens before serialization, chunking and padding in training, validation
and test loaders, for both HDF5 and legacy ragged inputs. Noise tagging marks
newly excluded clusters as class 1; finding treats them as non-target members,
negative in every eligible-track mask. Noise class-balance counts use the same
projected targets. The corrected labels change gradients while leaving the loss
formula unchanged. Generic source views, momentum/PID targets and original TPC
tasks retain their existing behavior.

The effective HDF5 contract records the policy and a minimum of four. Its noise
policy is \code{primary_unassigned_or_lt4_is_noise_v2}.
Existing running jobs retain their
frozen source and original targets. The noise checkpoint loader continues to
reject an unacknowledged data-contract mismatch; old checkpoint metadata is not
silently migrated.

Saved prediction scoring remains historical by default. Selecting
\code{--truth-policy primary_mc_min4_v1} explicitly projects each full event's
truth in memory and records the policy, implementation hash and relabelled-cluster
count. Predictions and saved input archives are unchanged. The metric profile
is still \code{mc_reconstruction_v1}; the truth policy is a separate required
comparison identity. Operational details and tests are documented in
\file{docs/evaluation/primary_mc_min4_v1.md}.

\section{Inputs, notation and units}

For an event $e$, let $V_e$ be the valid cluster indices after removing padding.
One ``hit'' in the metric names means one retained CVT \emph{cluster token}, not a
strip, ADC row, energy deposit, or generated particle. All clusters have unit
weight, regardless of their member count or energy.

Let $y_e(h)$ and $p_e(h)$ be the truth and final prediction labels. Define
\begin{align}
 T_{ei} &= \{h\in V_e:y_e(h)=i\}, & i&\ne b,\\
 S_e &= \bigcup_i T_{ei}, & U_e&=V_e\setminus S_e,\\
 P_{ej} &= \{h\in V_e:p_e(h)=j\text{ and prediction gate is true}\}, & j&\ge0.
\end{align}
Here $b=-1$ by default. $S_e$ is eligible truth support and $U_e$ is non-target
support. Truth track IDs and predicted track IDs are unrelated event-local
identifiers. Padding must be excluded by the validity mask; a negative truth
label other than $b$ on valid support is rejected by the new API.

Each valid cluster has at most one predicted label, so predicted sets are
disjoint. Negative predicted labels and a false prediction gate mean unassigned.
The scorer accepts final reconstruction assignments, not overlapping candidate
masks or raw query probabilities. Prediction thresholding is performed upstream
without truth assistance and is held fixed for comparisons.

Define $n_e$ and $m_e$ as the number of nonempty truth and prediction sets. For
each possible pair, write
\begin{align}
 a_{eij}&=|T_{ei}\cap P_{ej}|,& t_{ei}&=|T_{ei}|,& q_{ej}&=|P_{ej}|,\\
 c_{eij}&=\ratio{a_{eij}}{t_{ei}},& r_{eij}&=\ratio{a_{eij}}{q_{ej}},&
 J_{eij}&=\ratio{a_{eij}}{t_{ei}+q_{ej}-a_{eij}}.
\end{align}
$c$ is pair completeness, $r$ is pair purity and $J$ is intersection-over-union
(IoU). Their denominators contain \emph{all} members of the respective sets.
In particular, non-target clusters are not removed from prediction sizes.

\section{Two independent matching procedures}

\subsection{Dice assignment for the new reconstruction metrics}

The pair Dice score is
\begin{equation}
 D_{eij}=\frac{2a_{eij}}{t_{ei}+q_{ej}}
        =\frac{2c_{eij}r_{eij}}{c_{eij}+r_{eij}}
        =\frac{2J_{eij}}{1+J_{eij}}.
\end{equation}
The harmonic-mean expression is zero when the intersection is zero. Nonempty
sets make the first expression well defined. Scores lie in $[0,1]$.

Let $\mathcal M_e$ be a one-to-one set of pairs maximizing
\begin{equation}
 \max_{\mathcal M_e}\sum_{(i,j)\in\mathcal M_e}D_{eij}.
\end{equation}
The implementation uses SciPy's Hungarian assignment solver on negative Dice
weights. No IoU, purity, completeness or minimum-Dice acceptance cut is applied.
Only positive-overlap assigned pairs are retained. Solver pairs with zero
intersection are discarded and do not make a prediction ``MC-associated''.
Rows containing no target overlap are removed before solving; independent
non-target predictions therefore cannot influence assignment tie order.

Each truth and prediction can appear in at most one retained pair. A missed
truth track receives score zero. In a split, at most one fragment can supply its
truth track's Dice score. In a merge, at most one truth track can claim that
prediction. Other eligible track clusters still count in the prediction size.
Assignment is always event-local. Tracks from different events never compete.

\subsection{Assignment ties}

The primary total is the same for all exact maximizing assignments, but secondary
intersection and prediction-size totals can differ. The implementation orders
disjoint sets by their smallest original cluster index, then uses the solver's
deterministic result with the recorded SciPy version. It adds no epsilon weights
that could perturb the primary objective.

Provide stable, unique \code{point_ids} when comparing different serializations.
Without them, the input positions define the order. Relabeling track IDs does not
change the Dice assignment order. Exact tie-sensitive secondary results should
be reproduced with the same point identity and software versions. The legacy
matcher continues to use its original numerical label order for compatibility.

\subsection{Accepted matches for existing track efficiency}

Independently form the accepted-edge indicator
\begin{equation}
 A_{eij}=\mathbf1[J_{eij}\ge0.5\ \land\ r_{eij}\ge0.5\ \land\ c_{eij}\ge0.5].
\end{equation}
These are inclusive inequalities. The IoU cut remains mandatory. The legacy
solver maximizes the sum of weights
\begin{equation}
 W_{eij}=A_{eij}\bigl(\min(n_e,m_e)+1+J_{eij}\bigr),
\end{equation}
then retains accepted pairs only. The bonus first maximizes the number of
accepted matches and then their total IoU. Denote this result $\mathcal E_e$ and
its size $k_e$. It is not derived by thresholding $\mathcal M_e$.

The public \code{MatchConfig} continues to support explicitly configured legacy
cuts. Default Phase-1 scoring uses the values above; outputs record the config,
and accumulators reject different configs within a pooled result. No metric
change modifies the training loss's own Hungarian assignment.

\section{Primary and secondary reconstruction metrics}

Unless stated otherwise, sums below run over all events and their Dice pairs.
Define the sufficient totals
\begin{align}
 N_T&=\sum_e n_e,& N_S&=\sum_e|S_e|,\\
 Q_D&=\sum_e\sum_{(i,j)\in\mathcal M_e}D_{eij},&
 C&=\sum_e\sum_{(i,j)\in\mathcal M_e}a_{eij},\\
 H_A&=\sum_e\sum_{(i,j)\in\mathcal M_e}q_{ej}.&&
\end{align}
The associated support $R_e=\bigcup_{(i,j)\in\mathcal M_e}P_{ej}$ has
$\sum_e|R_e|=H_A$ because predictions are disjoint and matching is one-to-one.

\subsection{Primary truth-track mean Dice}
\begin{equation}
 \boxed{\metric{mc_track_dice_mean}=\frac{Q_D}{N_T}}
\end{equation}
Every supplied eligible truth track has one vote, including unrecovered tracks.
Events with many truth tracks contribute more votes than events with one track.
Empty-target events do not add primary denominator weight, but their diagnostic
counts and provenance still contribute where applicable.

\subsection{Hit completeness, purity and total impurity}
\begin{align}
 \metric{mc_hit_completeness_global}&=C/N_S,\\
 \metric{mc_hit_purity_global}&=C/H_A,\\
 \metric{mc_hit_impurity_global}&=(H_A-C)/H_A.
\end{align}
Completeness counts all eligible truth clusters in its denominator, including
clusters of missed tracks, unassigned clusters and fragments not claimed by the
one-to-one assignment. Purity uses only Dice-associated predictions in its
denominator, but every cluster in those predictions counts.

These are ratios of dataset totals, not means of pair purities or completenesses.
Their harmonic mean $2C/(N_S+H_A)$ generally differs from mean track Dice.
Maximizing total Dice also need not maximize $C$ by itself. All three quantities
describe the same Dice assignment; a second hit-count-maximizing association
would be a different metric.

Let $X$ count clusters from \emph{other eligible} tracks in associated predictions,
and let $U_A=\sum_e|R_e\cap U_e|$. Then
\begin{equation}
 H_A=C+X+U_A.
\end{equation}
The additional diagnostics are
\begin{align}
 \metric{mc_other_track_contamination_global}&=X/H_A,\\
 \metric{mc_nontarget_contamination_global}&=U_A/H_A.
\end{align}
Their sum is total impurity. Excluded MC and overlay are both in $U_A$ if they
contaminate an associated prediction. Their independent predictions contribute
to neither $H_A$ nor $C$.

\subsection{Verified overlay contamination and leakage}

Let $B_e\subseteq U_e$ be clusters verified to consist entirely of measured
overlay members. Define
\begin{equation}
 H_B=\sum_e|B_e|,\qquad B_A=\sum_e|R_e\cap B_e|.
\end{equation}
With adequate provenance on the relevant support,
\begin{align}
 \metric{mc_bg_contamination_global}&=B_A/H_A,\\
 \metric{bg_leakage_to_mc_global}&=B_A/H_B.
\end{align}
The first asks how much of the associated prediction support is verified pure
overlay. The second asks how much of the dataset's verified pure overlay entered
those associated predictions. It includes overlay in empty-target events in
$H_B$, so the latter statistic depends on event-cohort composition.

Both names refer to the pure-overlay cluster definition in this version. They do
not count every cluster containing even one overlay member, and do not estimate
the fraction of member strips from overlay. Mixed-cluster fractions are reported
separately. A mixed eligible cluster is still a whole eligible cluster for the
track score under the existing ownership rule; provenance does not fractionally
reweight or relabel it.

When some associated clusters have \code{UNKNOWN} provenance, contamination is
\code{null}. When any valid cluster in the dataset has \code{UNKNOWN} provenance,
leakage is \code{null}, because its overlay denominator cannot be fully classified.
The separate \metric{mc_bg_contamination_lower_bound_global} always reports
$B_A/H_A$ when $H_A>0$. Unknown clusters are never counted as verified overlay.
This lower bound can be zero without establishing that contamination is zero.

\subsection{Accepted-match efficiency and signal ARI}
\begin{equation}
 \metric{track_efficiency_global}=\frac{\sum_e k_e}{N_T}.
\end{equation}
This retained metric uses $\mathcal E_e$, not $\mathcal M_e$.
\metric{ari_signal} retains its event-macro legacy definition, described below.
Neither diagnostic becomes a new checkpoint-selection rule in Phase 1.

\section{Provenance representation and accounting}

\begin{longtable}{p{0.06\linewidth}p{0.31\linewidth}p{0.55\linewidth}}
\toprule Code & Category & Established information\\\midrule\endhead
0 & \code{UNKNOWN} & Composition is unavailable or cannot be classified by the supplied evidence.\\
1 & \code{MC_ONLY} & All observed members are known MC; track eligibility is a separate label decision.\\
2 & \code{OVERLAY_ONLY} & All observed members are verified measured overlay.\\
3 & \code{MIXED} & Both MC and measured-overlay members are explicitly established.\\
4 & \code{MC_PRESENT_UNRESOLVED} & Some MC is established; whether remaining members are other MC or overlay is unresolved. The cluster cannot be pure overlay.\\
\bottomrule
\end{longtable}

These categories partition valid clusters. \code{MC_ONLY} on a non-target cluster
means MC outside the supplied eligible support; the evaluator does not infer its
particle identity or the reason for exclusion. \code{MC_PRESENT_UNRESOLVED} must
not be described as verified mixed or pure excluded MC.

\subsection{Conservative mapping of existing HDF5 fields}

\code{cluster_provenance_from_hdf5} consumes aligned cluster arrays:
\begin{itemize}
\item \code{background_only=true} establishes \code{OVERLAY_ONLY}.
\item If supplied, \code{ownership=0} must agree exactly with that flag.
      Ownership 1 is strict primary majority, 2 contested, and 3 other nonprimary.
\item \code{origin_noise=false} establishes some MC contribution, hence
      \code{MC_PRESENT_UNRESOLVED}; it does not establish pure MC composition.
\item Ownership 1 or 2 also establishes primary MC evidence. Contradictory
      \code{origin_noise=true} is rejected.
\item Ownership 1 with direct-primary \code{purity=1} establishes \code{MC_ONLY}.
      The comparison is exact and uses no tolerance that could promote a
      nearly-pure cluster to pure MC.
\item Remaining cases are \code{UNKNOWN}. In particular, no simulated origin
      contributor is not sufficient evidence of measured overlay.
\end{itemize}

The compact HDF5 fields cannot, in general, distinguish all-MC mixtures from
MC/overlay mixtures. The helper therefore never invents \code{MIXED}. Callers
with an existing, explicit member-provenance sidecar may supply that category
directly. No producer change or new truth-track inference is required.

Supplying pure-overlay provenance for a cluster labelled eligible MC is a
contradiction and raises an error. Missing provenance defaults to
\code{UNKNOWN}; primary Dice, completeness, total purity, efficiency and legacy
metrics remain computable.

\subsection{Counts, coverage and exact decompositions}

For each lower-case category name $g$, the output includes
\begin{itemize}
\item \code{n_<g>_points}: category count on all valid support;
\item \code{n_mc_associated_<g>_points}: category count on $R_e$;
\item \code{n_mc_nontarget_<g>_points}: category count on $R_e\cap U_e$.
\end{itemize}
Let $Z_A$ be the associated \code{UNKNOWN} count and $Z$ the dataset unknown count.
Then
\begin{align}
 \metric{mc_provenance_coverage_global}&=(H_A-Z_A)/H_A,\\
 \metric{provenance_coverage_global}&=(|V|-Z)/|V|.
\end{align}
This is coverage of the categorical classification, not the fraction of clusters
whose every member origin is known. Unresolved MC presence is a classified
category, and its fraction remains visible separately.

The \code{mc_associated_<g>_fraction_global} family divides associated category
counts by $H_A$; the five values sum to one when $H_A>0$.
The \code{mc_nontarget_<g>_fraction_global} family divides non-target associated
category counts by $H_A$; its sum is non-target contamination. Thus
\begin{equation}
 \frac{C}{H_A}+\frac{X}{H_A}+\sum_g\frac{U_{A,g}}{H_A}=1.
\end{equation}
Associated composition includes correct and wrong eligible clusters as well as
non-target clusters. It must not be added to the impurity decomposition again.

\subsection{Exported sufficient counts}

These fields are sums of event counts in a dataset summary. They are not
averages; \code{mc_track_dice_sum} is the one score total in this table.
\begin{longtable}{p{0.54\linewidth}p{0.39\linewidth}}
\toprule Output field & Meaning\\\midrule\endhead
\metric{n_events} & Number of complete events scored.\\
\metric{n_points} & Total valid clusters, $|V|$.\\
\metric{n_signal_points} & Eligible truth clusters, $N_S$.\\
\metric{n_background_points} & All non-target clusters, $|U|$.\\
\metric{n_true_tracks} & Supplied eligible truth tracks, $N_T$.\\
\metric{n_pred_tracks} & Nonempty final predictions, $\sum m_e$.\\
\metric{n_matched_tracks} & Accepted matches, $\sum k_e$.\\
\metric{n_mc_dice_matches} & Positive-overlap Dice associations.\\
\metric{mc_track_dice_sum} & Sum of Dice scores, $Q_D$.\\
\metric{n_mc_correct_points} & Assigned pair intersections, $C$.\\
\metric{n_mc_associated_pred_points} & Associated prediction support, $H_A$.\\
\metric{n_mc_other_track_points} & Other eligible clusters in that support, $X$.\\
\metric{n_mc_nontarget_points} & Non-target clusters in that support, $U_A$.\\
\metric{n_split_true_tracks} & Truth sets with more than one positive prediction overlap.\\
\metric{n_merged_pred_tracks} & Predictions with more than one positive truth overlap.\\
\metric{n_non_target_only_pred_tracks} & Predictions with no eligible truth overlap.\\
\metric{n_mc_overlap_unassociated_pred_tracks} & Positive-overlap predictions not in the Dice assignment.\\
\metric{n_background_assigned_points} & Non-target clusters assigned to any prediction.\\
\metric{n_signal_unassigned_points} & Eligible clusters left unassigned.\\
\metric{n_truth_tracks_below_policy_min_clusters} & Supplied truth sets with fewer than four clusters; audit only.\\
\metric{n_events_with_provenance} & Events supplied a provenance array, even if some or all entries are unknown.\\
\bottomrule
\end{longtable}

\section{Aggregation and undefined values}

\subsection{Dataset-global ratios}

For a per-event ratio $r_e=x_e/d_e$, its global variant is
\begin{equation}
 r_{\mathrm{global}}=\frac{\sum_e x_e}{\sum_e d_e}.
\end{equation}
When all defined denominators are positive, this is a denominator-weighted event
average. The weight is truth tracks for mean Dice, truth clusters for global
completeness, associated predicted clusters for global purity, and verified pure
overlay clusters for leakage. These are different weighting choices.

Accumulators sum integer counts and float64 score totals across complete events.
They may merge batch accumulators. They never average batch means. Integer totals
are exactly batch invariant; floating sums may differ at roundoff level due to
addition order. Tests require agreement to $10^{-12}$ on fixed predictions.
This guarantee does not imply that changing a GPU inference batch size produces
bit-identical predictions.

\subsection{Event-macro diagnostics}

For a metric defined on event subset $F$, the macro mean is
\begin{equation}
 r_{\mathrm{event\ macro}}=\frac{1}{|F|}\sum_{e\in F}r_e.
\end{equation}
The explicit primary companion is \metric{mc_track_dice_event_macro}. New ratio
families have \code{_event_macro} companions. Historical macro fields retain
their old unsuffixed names and also have explicit aliases. The new macro
observation counts are recorded in \code{mc_event_macro_observations}.

Macro overlay values exclude events where required provenance is unknown. A
pooled global value can therefore be undefined even when an available-event
macro value is defined. Always report the coverage and counts; these are not
interchangeable populations.

\subsection{Zero denominators}

All undefined ratios are Python \code{None}/JSON \code{null}, never NaN, infinity,
or an invented zero or one. Important cases are:
\begin{longtable}{p{0.43\linewidth}p{0.50\linewidth}}
\toprule Condition & Result\\\midrule\endhead
No eligible truth tracks & Mean Dice, track efficiency and hit completeness undefined.\\
Truth exists, all tracks missed & Mean Dice and hit completeness zero.\\
No positive-overlap Dice associations & Associated purity, impurity, contamination and coverage undefined.\\
No verified pure overlay, complete provenance & Overlay contamination zero if associated support exists; leakage undefined.\\
Unknown associated provenance & Overlay contamination undefined; its separately named lower bound remains available.\\
Unknown provenance anywhere & Overlay leakage undefined.\\
Empty input dataset & Zero counts and undefined ratios; no empty-dataset perfect score.\\
\bottomrule
\end{longtable}

\section{Historical metrics retained exactly}

The original event scorer \code{event_track_metrics} retains its definition.
So do \code{summarize_event_metrics} and the streaming accumulator
\code{EventMetricAccumulator}. Their accepted matcher is the current version-2
eligible-first matcher. Older reports made with a different matcher must retain
their own version/provenance; they must not be relabelled as equivalent.

For this section, let $F_e$ count truth tracks overlapping more than one
prediction, and $G_e$ count predictions overlapping more than one truth track.
Any positive cluster overlap counts; there is no significant-overlap cut.
Let $L_e$ be eligible clusters left unassigned, $U^+_e$ non-target clusters
assigned to any prediction, and $H^+_e=\sum_jq_{ej}$ all assigned clusters.

\begin{longtable}{p{0.40\linewidth}p{0.53\linewidth}}
\toprule Historical field & Event definition and interpretation\\\midrule\endhead
\metric{track_efficiency} & $k_e/n_e$: fraction of truth tracks with accepted matches.\\
\metric{track_purity} & $k_e/m_e$: fraction of predictions with accepted matches; not hit purity.\\
\metric{matched_iou_mean} & Mean $J_{eij}$ over $\mathcal E_e$.\\
\metric{matched_purity_mean} & Mean $r_{eij}$ over $\mathcal E_e$.\\
\metric{matched_efficiency_mean} & Mean $c_{eij}$ over $\mathcal E_e$.\\
\metric{fake_rate} & $(m_e-k_e)/m_e$. Historical name: unmatched predictions are not established physical fakes.\\
\metric{miss_rate} & $(n_e-k_e)/n_e$.\\
\metric{split_rate} & $F_e/n_e$, using positive target overlaps.\\
\metric{merge_rate} & $G_e/m_e$, using positive target overlaps.\\
\metric{background_rejection} & $(|U_e|-U^+_e)/|U_e|$: non-target clusters left unassigned. Not overlay-only rejection.\\
\metric{background_contamination} & $U^+_e/H^+_e$: non-target clusters in all predictions, including independent non-target predictions.\\
\metric{signal_loss_to_background} & $L_e/|S_e|$: eligible clusters left unassigned. Does not count wrong-track assignment.\\
\metric{ari_signal} & Adjusted Rand index on eligible support, with the legacy degenerate-case convention below.\\
\metric{ari_with_background} & Adjusted Rand index on valid support, treating all non-target truth as one class.\\
\bottomrule
\end{longtable}

Dataset values of these unsuffixed fields are means across events where the
quantity is defined. In particular, a historical matched-quality mean first
averages accepted pairs within an event, then averages defined events equally.
It is not a dataset-pooled accepted-track average.

The existing \metric{track_efficiency_global} and
\metric{track_purity_global} are $\sum k_e/\sum n_e$ and $\sum k_e/\sum m_e$.
They remain unchanged. The former shares the primary truth-track denominator
but has a different matching objective and thresholded numerator.

\subsection{ARI equations and edge conventions}

On the chosen support, let $n_{ij}$ be the truth/prediction contingency table,
$a_i=\sum_j n_{ij}$, $b_j=\sum_i n_{ij}$, and $n=\sum_{ij}n_{ij}$. Write
\begin{equation}
 I=\sum_{ij}\binom{n_{ij}}2,\quad A=\sum_i\binom{a_i}2,\quad
 B=\sum_j\binom{b_j}2,\quad Q=\binom n2.
\end{equation}
The adjusted Rand index is
\begin{equation}
 \mathrm{ARI}=\frac{I-AB/Q}{(A+B)/2-AB/Q},
\end{equation}
with the library's conventions for degenerate partitions. The legacy wrapper
uses sklearn's ARI on eligible support only when there are more than one
eligible cluster and more than one truth track. Otherwise, with any eligible
support, it returns 1 if that support has at most one predicted label and 0 if
it has multiple predicted labels. With no eligible support it returns undefined.

Consequently, an entirely missed single-track event can have
\metric{ari_signal}=1 because all its clusters share the unassigned label.
Signal ARI also ignores contaminating non-target clusters. These are reasons to
retain ARI as a partition diagnostic rather than interpreting it as complete
reconstruction quality. No historical convention is silently corrected here.

Inclusive ARI is computed only when there are more than one valid cluster and
more than one distinct truth label on valid support; otherwise it is undefined.
All unassigned predictions share label $-1$. Grouping all non-target truth into
one class can penalize independent non-target reconstructions, so inclusive ARI
does not express the primary benchmark objective.

\subsection{Additional global diagnostics and aliases}

New global variants divide summed numerators by summed denominators from the
table. These cover split, merge and miss rates, background rejection, background
contamination and signal loss to background. Each uses its existing field name
with the explicit \code{_global} suffix.
\code{matched_<quality>_mean_global} pools all
accepted pairs across the dataset, for qualities IoU, purity and efficiency.

\metric{unmatched_predicted_fraction_global} is
$(\sum m_e-\sum k_e)/\sum m_e$; its macro companion has exactly the historical
\code{fake_rate} meaning without asserting physical fakeness. No old field is
removed or silently renamed.

\metric{non_target_only_predicted_fraction_global} counts predictions with
\emph{zero} eligible truth overlap divided by all predictions. It does not
establish that these predictions are valid tracks or fakes.
\metric{mc_overlap_unassociated_predicted_fraction_global} counts predictions
with positive eligible overlap that were not selected by the Dice assignment,
divided by all predictions. These can include split fragments or competing mixed
predictions. Their macro companions use the corresponding per-event ratios.

No dataset-global ARI is made by concatenating event-local IDs; doing so would
identify unrelated track labels across events. The preserved ARI aggregation
is explicitly event macro.

\section{Worked examples}

All examples use final disjoint assignments and unit-weight clusters. Unless
stated otherwise, there is complete MC-only provenance for target clusters.

\subsection{Perfect reconstruction plus independent non-target tracks}
Two truth tracks each have four clusters. Predictions recover each exactly and
also form a separate four-cluster overlay-only track. Mean Dice, completeness,
purity and track efficiency are all 1. Overlay contamination and leakage to the
associated predictions are 0. The historical prediction purity is $2/3$ and
the historical fake rate is $1/3$. That disagreement is intentional: the extra
prediction's physical validity is outside the primary truth contract.

\subsection{One four-cluster track split into two equal fragments}
Each fragment has two correct clusters. One can be associated; its Dice is
$2(2)/(4+2)=2/3$. Mean Dice is $2/3$, completeness is $1/2$, purity is 1, and
split rate is 1. The legacy matcher accepts one fragment exactly at IoU and
completeness $0.5$, so track efficiency is 1. Continuous quality and thresholded
efficiency reveal different aspects of this reconstruction.

\subsection{Two four-cluster tracks merged into one prediction}
The eight-cluster prediction can be associated to only one truth track. Its
Dice is $2(4)/(4+8)=2/3$, while the other truth track contributes zero. Mean Dice
is therefore $1/3$, completeness and purity are each $1/2$, and merge rate is 1.
The other truth track's four clusters contribute to other-track impurity.
One legacy accepted match gives track efficiency $1/2$.

\subsection{Different impurity sources}
Start with two perfectly recovered four-cluster targets. Add four non-target
clusters to the first prediction: one MC-only, one pure overlay, one explicitly
mixed, and one MC-present unresolved. A second pure-overlay cluster forms an
independent prediction. Then
\begin{align}
 N_T&=2,&N_S&=8,&C&=8,&H_A&=12,&H_B&=2,&B_A&=1,\\
 \overline D&=\frac{2/3+1}{2}=\frac56,&
 \text{completeness}&=1,&\text{purity}&=\frac23.
\end{align}
Total impurity is $4/12$, but verified pure-overlay contamination is $1/12$ and
leakage is $1/2$. Each of the four non-target provenance categories contributes
$1/12$ to the impurity decomposition. None is promoted to a truth target.

\subsection{Truth-track mean versus event macro versus hit completeness}
One event has a perfectly recovered four-cluster track. Another has three
eight-cluster tracks, all missed. Mean Dice is $1/4$, event-macro Dice is $1/2$,
and global hit completeness is $4/(4+24)=1/7$. Averaging the two event scores or
the two batch scores would produce the wrong primary result.

\subsection{Dice and accepted matching can disagree}
Take two truth tracks with sizes 12 and 4, and predictions with sizes 12 and 4,
whose intersection table is
\begin{equation}
 \begin{pmatrix}8&4\\4&0\end{pmatrix}_{\text{prediction}\times\text{truth}}.
\end{equation}
The Dice matrix is
$\left(\begin{smallmatrix}2/3&1/2\\1/2&0\end{smallmatrix}\right)$.
Maximum-total-Dice matching selects the off-diagonal pairs for total 1 and mean
$1/2$. Only the upper-left edge meets the IoU threshold; the independent legacy
matcher accepts that edge. Thresholding the Dice assignment would incorrectly
report zero accepted matches.

\section{Implementation and saved-prediction contract}

\subsection{Shared API}
The implementation is self-contained in
\file{train/downstream/track_finding_metrics.py}, retaining its NumPy, SciPy and
sklearn dependencies. The opt-in interface is:
\begin{Code}
row = event_mc_reconstruction_metrics(
    truth, pred, valid_mask=valid, pred_signal_mask=gate,
    provenance=provenance, point_ids=original_cluster_index,
)
accumulator = MCReconstructionAccumulator()
accumulator.add(row)          # one complete event
accumulator.merge(other)      # optional disjoint batch/shard
summary = accumulator.summary()
\end{Code}
An event row contains every historical field, unchanged, plus
\code{dice_matches}, the new counts, \code{mc_track_dice_sum} and a
\code{mc_event_metrics} dictionary. The historical \code{matches} list remains
the accepted-match list. Each Dice pair records truth ID, prediction ID,
intersection, both sizes and Dice. One-to-one uniqueness and positive overlap
can be audited directly.

The summary includes old outputs, explicit macro aliases, new global metrics,
new macro metrics, counts, profile version and accepted-match configuration.
It does not change \code{FINDING_METRICS_VERSION=2}, which belongs to the legacy
training/checkpoint contract. Mixing legacy-only rows into the new accumulator
raises an error rather than fabricating missing counts.

\subsection{Offline NPZ scorer}
Run from the repository root:
\begin{Code}
python -m train.downstream.eval.rescore_track_finding \
  --input /path/to/predictions.npz \
  --prediction-key pred coatjava \
  --output /path/to/new_summary.json \
  --events-output /path/to/new_events.jsonl
\end{Code}
The archive must contain:
\begin{itemize}
\item \code{offsets}: integer array of length events plus one, beginning at zero,
      nondecreasing, ending at the flat point count;
\item \code{truth}: flat integer supplied target labels;
\item one flat final-prediction label array per \code{--prediction-key};
\item optional aligned \code{point_ids}, and either provenance codes or raw
      \code{background_only}, \code{ownership}, \code{origin_noise}, \code{purity};
\item optional event arrays \code{source_event_index}, \code{source_id},
      \code{original_ordinal}; duplicate available identities are rejected.
\end{itemize}
Stored labels must already incorporate the native assignment gate and threshold.
They must not contain padding or duplicated overlapping event chunks. Empty
events may have repeated consecutive offsets. Pickled/object arrays are not used.

\code{--events-per-batch} controls CPU aggregation batches only. It cannot change
the saved predictions. Output files must be new; historical reports are never
overwritten. A successful summary contains input and scorer hashes, dependency
versions, profile/config, provenance source and the full-event verification flag.
The optional event JSONL is a streaming diagnostic; if scoring fails, it may be
partial and is not a completed summary.

Without an HDF5 join, the caller is responsible for full-event support and
truth/provenance alignment. The receipt explicitly records that full-event
membership was not independently verified. An old archive lacking provenance
can still yield the primary score and non-overlay metrics; it cannot establish
overlay contamination from $-1$ truth labels alone.

\subsection{Joining existing HDF5 provenance safely}
Supply all of:
\begin{Code}
  --hdf5-root /path/to/published/dataset \
  --split validation --plan-id EXACT_PUBLISHED_PLAN_ID
\end{Code}
The archive must also contain all three event identity arrays and original
\code{point_ids}. For every event the scorer checks the published plan ID,
physical source ID and original ordinal, and a complete permutation of the HDF5
cluster indices. By default, stored truth must exactly match the corresponding
HDF5 \code{primary_segment}. With the explicit minimum-four policy, stored truth
may equal either the original labels or their exact projected labels. The scorer
then applies that projection and permutes the existing provenance into prediction
order. Other mismatches fail; no coordinate-based alignment or arbitrary truth
repair is attempted.

This path deliberately refuses archives whose HDF5 origin or cluster ordering
cannot be proved. Existing checkpoints remain usable, but inference/export
must provide those identities before overlay-aware rescoring can be trusted.
The evaluator does not add provenance to model inputs. The separate loader patch
enforces eligibility as described in Section~\ref{sec:cut}.

\section{Validation requirements and evidence}

The CPU-only test module is
\file{train/downstream/tests/test_mc_reconstruction_metrics.py}. It checks:
\begin{enumerate}
\item Perfect targets, independent non-target predictions, and retained
      unmatched-prediction diagnostics.
\item Splits, merges, fully missed tracks and exact legacy threshold boundaries.
\item Different Dice and accepted assignments, including the counterexample above.
\item Exhaustive small-matrix assignment oracles, positive intersections and
      one-to-one uniqueness.
\item Truth-track weighting versus event macro and cluster weighting.
\item Full impurity/provenance decompositions, excluded MC, verified overlay,
      explicitly mixed clusters and missing provenance.
\item Empty datasets/events, zero denominators, padding, prediction gates,
      malformed arrays, immutable labels and short-track auditing.
\item Reproducible ties under label renaming and point serialization with
      stable original indices.
\item Batch sizes 1, 2, 7, 16, 64 and 100, with reverse-order accumulator merging.
\item Exact per-event legacy-field parity and numerical summary parity.
\item Saved-archive round trips and HDF5 event/point identity join failures.
\end{enumerate}

Run the focused suite with
\begin{Code}
python -m unittest \
  train.downstream.tests.test_mc_reconstruction_metrics -v
\end{Code}
Existing metric, validation-history and HDF5 downstream tests additionally
exercise compatibility. Some campaign tests require
\code{PYTHONPATH=.:train/downstream/campaign} under the repository's existing
import layout. A source-AST comparison verifies that pre-existing metric
functions/classes are unchanged, rather than relying only on shared-code parity.

A bounded real-data validation uses the locally available qualification HDF5
pilot, saved native COATJAVA assignments, deliberately permuted original cluster
indices, and CPU scoring batches 1 and 64. It checks the full provenance join and
pooled metric agreement. This is evaluator validation, not a new model ranking,
campaign test result or confirmation that every input satisfies eligibility.
Dated evidence belongs under
\file{research_reports/mc_reconstruction_phase1_20261110/}; it is not the normative
definition of the metrics.

\section{Versioning and future integration boundary}

The profile identifier \code{mc_reconstruction_v1} fixes the truth-label contract,
Dice objective, association rule, tie convention, weighting, provenance mapping,
undefined-value rules and field meanings in this specification. A scientific
definition change requires a new profile version. An implementation correction
must be accompanied by tests and a recorded source hash; prior artifacts retain
their original receipts. Changing aggregation requires a distinctly named field.

The Phase-1 code does not call the new metric from training, change the model
selection callback, reschedule jobs, rewrite checkpoints, or rank all saved
checkpoints. Later integration should use this same tested API, preserve fixed
validation event identities and final assignment settings, and export the
original point/provenance identities needed for reproducibility. Any later
checkpoint-selection change must be explicit and use validation data; final
test results remain a separate evaluation.

The intended four-cluster policy remains fixed. If an input's short-track audit
is nonzero, explain the input-contract discrepancy before claiming a
policy-compliant benchmark result. This is an input audit, not permission to
silently change the scientific target population.

\section{Implementation sources and reproducible document build}

Normative implementation entry points:
\begin{itemize}
\item \file{train/downstream/track_finding_metrics.py}: legacy metrics, new
      matching, provenance conversion and accumulators.
\item \file{train/downstream/eval/rescore_track_finding.py}: saved-label scorer,
      optional HDF5 join and receipt.
\item \file{fm4npp/datasets/cvt_downstream.py}: HDF5 projection, with the authorized
      cut for adapter supervision; raw views preserve stored labels.
\item \file{fm4npp/datasets/track_eligibility.py}: shared full-event cut, version
      and legacy ragged label views.
\item \file{fm4npp/datasets/cvt_hdf5.py}: existing published-manifest reader.
\item \file{train/downstream/track_finding_contract.py}: retains the legacy
      metric version and checkpoint guards.
\end{itemize}

The provenance mapping was checked against the existing
\code{cvt-clusters-h5-v1} schema and producer functions
\code{strict_ownership}, \code{clusters}, and \code{compact_primary_summary} in
the local \code{cvt-cluster-skim} package. Persisted fields, not inferred particle
identities, inform the new diagnostic categories.

This LaTeX file is the versioned source of the accompanying PDF. Build from the
repository root using
\begin{Code}
mkdir -p /tmp/mc-reconstruction-spec
latexmk -pdf -interaction=nonstopmode -halt-on-error \
  -outdir=/tmp/mc-reconstruction-spec \
  docs/evaluation/mc_reconstruction_v1.tex
\end{Code}
Render the resulting PDF to images for visual verification before publishing a
revised document. Equations and explanatory text must remain consistent with
the executable tests.

\end{document}
